\section{Análisis de la encuesta a docentes}

La encuesta fue realizada por cinco docentes del programa identificados como D1 a D5. El instrumento segmentó los bloques temáticos, permitiendo que los cinco diligenciaran el bloque de percepción general; D1 y D5 respondieron lo relacionado a la programación imperativa y programación orientada a objetos, teniendo en cuenta que dictan en primero y segundo semestre. Por otro lado, D2, D3 y D4 realizaron el bloque de programación lógica y paradigmas avanzados, lo que corresponde a quinto semestre en adelante. Si bien, el perfil directamente relacionado con FLP es el de los docentes de semestres avanzados, las respuestas de D1 y D5 permiten caracterizar las bases con las que los estudiantes llegan a la asignatura, lo que es un factor determinante en su capacidad de asimilar contenidos de alta abstracción \cite{zagami2023}.

\subsection{Percepción general de los estudiantes}

Este bloque obtuvo la respuesta de los cinco docentes y permite constuir un diagnóstico transversal del programa. La Tabla \ref{tab:percepcion_docentes} presenta los resultados, agrupando primero a los docentes de semestres avanzados (D2, D3, D4) y luego a los de primeros semestres (D1, D5), para facilitar la comparación por grupo.

\begin{table}[H]
\centering
\caption{Percepción general de los docentes sobre habilidades de los estudiantes ($n=5$)}
\label{tab:percepcion_docentes}
\footnotesize
\setlength{\tabcolsep}{3pt}

\begin{tabular}{
>{\raggedright\arraybackslash}p{5.0cm}
>{\centering\arraybackslash}p{1.8cm}
>{\centering\arraybackslash}p{1.8cm}
>{\centering\arraybackslash}p{1.8cm}
>{\centering\arraybackslash}p{1.8cm}
>{\centering\arraybackslash}p{1.8cm}
}
\hline
\textbf{Afirmación} & \textbf{D2 (5°+)} & \textbf{D3 (5°+)} & \textbf{D4 (5°+)} & \textbf{D1 (1°-2°)} & \textbf{D5 (1°-2°)} \\
\hline

Demuestran habilidades básicas de pensamiento lógico &
De acuerdo &
Tot. acuerdo &
Tot. acuerdo &
Neutral &
En desacuerdo \\

Tienen dificultades para analizar un problema antes de programar &
En desacuerdo &
De acuerdo &
De acuerdo &
En desacuerdo &
Neutral \\

Logran depurar errores de manera efectiva &
De acuerdo &
De acuerdo &
Neutral &
En desacuerdo &
Neutral \\

Pueden dividir problemas complejos en partes más pequeñas &
Tot. acuerdo &
Neutral &
En desacuerdo &
En desacuerdo &
Neutral \\

Comprenden la importancia de la abstracción &
Tot. acuerdo &
Tot. acuerdo &
En desacuerdo &
En desacuerdo &
Neutral \\
\hline

\end{tabular}

\vskip 4pt
\footnotesize\textit{Fuente: elaboración propia.}
\end{table}

\FloatBarrier

La lectura de la Tabla \ref{tab:percepcion_docentes} revela una brecha de percepción entre grupos. Los docentes de semestres avanzados tienden a valorar positivitamente las habilidades de sus estudiantes que los de primeros semestres, lo cual es un resultado esperable dado que los estudiantes de semestres superiores han superado los filtros académicos iniciales. D1 se muestra consistentemente en desacuerdo en cuatro de los cinco ítems, mientras que D5 se mantiene neutral, ambos describiendo un perfil de estudiante de primeros semestre con habilidades lógicas y de abstracción mínimas. 

Además, se observa una varidad de respuestas dentro del grupo avanzado, mientras D2 y D3 coinciden en que los estudiantes de semestre superiores comprenden la importancia de la abstracción, D4 está en desacuerdo con esta afirmación. Lo mismo ocurre con la capacidad de dividir problemas complejas, donde D2 responde "Totalmente de acuerdo" y D4 "En desacuerdo". Estas diferencias dentros del mismo grupo de docentes indica que las habilidades de abstracción y descomposición de problemas no están consolidadas de manera uniforme en los estudiante. 

Esta observación es coherente con lo expuesto por Zagami, quien señala que los estudiantes enfrentan una alta demanda cognitiva al intentar procesar conceptos abstractos simultáneamente con habilidades prácticas; cuando la capacidad de la memoria de trabajo se ve excedida por esta multiplicidad de elementos, se generan bloqueos y sentimientos de frustración en el aprendiz \cite{zagami2023}. Asimismo, se reconoce que la naturaleza intrínsecamente abstracta de los algoritmos y construcciones teóricas en esta disciplina representa un desafío que dificulta la comprensión intuitiva \cite{Stamenkovic2024AWE}.

\subsection{Dificultades con paradigmas avanzados}

Los docentes D2, D3 y D4 respondieron el bloque sobre los paradigmas directamente relacionados con los contenidos de FLP. La Tabla \ref{tabla:percepcion_paradigmas_docentes} muestra sus respuestas, evidenciando una percepción más positiva en D2 y D3, mientras que D4 se muestra más crítico, especialmente en relación con la capacidad de los estudiantes para seleccionar el paradigma adecuado para resolver un problema.

\begin{table}[H]
\centering
\caption{Percepción de los docentes de semestres avanzados sobre programación lógica y paradigmas avanzados ($n=3$)}
\label{tabla:percepcion_paradigmas_docentes}
\footnotesize
\setlength{\tabcolsep}{3pt}

\begin{tabular}{
>{\raggedright\arraybackslash}p{7cm}
>{\centering\arraybackslash}p{2.5cm}
>{\centering\arraybackslash}p{2.5cm}
>{\centering\arraybackslash}p{2.5cm}
}
\hline
\textbf{Afirmación} & \textbf{D2 (5°+)} & \textbf{D3 (5°+)} & \textbf{D4 (5°+)} \\
\hline

Comprenden los conceptos fundamentales de programación lógica &
Tot. acuerdo &
Tot. acuerdo &
Tot. acuerdo \\

Tienen dificultades con el razonamiento lógico aplicado a la programación &
De acuerdo &
De acuerdo &
En desacuerdo \\

Logran modelar problemas usando reglas lógicas &
Tot. acuerdo &
Neutral &
De acuerdo \\

Comprenden la diferencia entre los distintos paradigmas &
Tot. acuerdo &
De acuerdo &
En desacuerdo \\

Pueden seleccionar el paradigma adecuado para resolver un problema &
De acuerdo &
De acuerdo &
Tot. en desacuerdo \\
\hline

\end{tabular}

\vskip 4pt
\footnotesize\textit{Fuente: elaboración propia.}
\end{table}

\FloatBarrier

El único ítem en el que los tres docentes coinciden es la comprensión de los conceptos fundamentales de programación lógica, donde todos respondieron "Totalmente de acuerdo". Sin embargo, esta percepción contrasta con el resto de las respuestas lo que sugiere que, saber qué es la programación lógica no implica saber aplicarla para resolver problemas.

Las opiniones también difieren respecto al razonamiento lógico aplicado a la programación. D2 y D3 consideran que los estudiantes presentan dificultades en este aspecto, mientras que D4 no lo percibe así. Estas diferencias pueden estar relacionadas con el contexto y el tipo de actividades realizados, ya que en cursos más guiados, algunas dificultades pueden pasar desapercibidas, pero, en asignaturas como FLP, donde los estudiantes deben construir soluciones con mayor autonomía, estas limitaciones resultan  más evidentes.

Así mismo, la capacidad de comprender las diferencias entre paradigmas y seleccionar el adecuado para un problema concentra la mayor dispersión del bloque, debido a que D2 y D3 tienen una percepción positiva,  pero, D4 está respectivamente ``En desacuerdo'' y ``Totalmente en desacuerdo''. Esta percepción negativa demuestra el salto cognitivo que representa abandonar el paradigma imperativo para adoptar el paradigma funcional \cite{nemeth2019}, dificultad que se agrava cuando el lenguaje utilizado impone una carga sintáctica adicional. 